\documentclass[aspectratio=169,11pt]{beamer}

\usetheme{default}
\usecolortheme{dove}
\setbeamertemplate{navigation symbols}{}
\setbeamertemplate{footline}[frame number]
\setbeamerfont{frametitle}{size=\large}
\setbeamertemplate{itemize item}{\textbullet}

\usepackage{amsmath}
\usepackage{amssymb}
\usepackage{booktabs}

\newcommand{\uloc}{\hat{u}}
\newcommand{\uext}{\hat{v}}
\newcommand{\x}{\mathbf{x}}
\newcommand{\dOmega}{\,\mathrm{d}\Omega}

\title{Pedestrian exposure in a city square}
\subtitle{The method in five equations}
\author{Robin Wydaeghe}
\date{}

\begin{document}

\begin{frame}
  \titlepage
\end{frame}

\begin{frame}{The quantity}
  A pedestrian in a city square receives power from many base stations at once.
  The buildings block part of it and reflect part of it back. That is one number:
  \begin{equation*}
    \boxed{\;\chi = \frac{S}{S_0}\;}
  \end{equation*}
  \begin{itemize}
    \item $S$ is the power density arriving at the head in the real square, in W/m$^2$.
    \item $S_0$ is what the same network delivers at the same point over open ground.
    \item $\chi$ is dimensionless. It is $1$ in free space, below $1$ where the
      square shadows more than it reflects, above $1$ where it reflects more than
      it shadows.
  \end{itemize}
  \vspace{0.4em}
  A ratio rather than an absolute level, because the transmit power, the site
  count and the site positions are not public at millimetre wave and are not
  properties of the square. Dividing removes all three.
\end{frame}

\begin{frame}{Notation}
  Unit vectors carry a hat. Elevation $\alpha$ is measured up from the horizon,
  so $\alpha = 0$ is level with the head.
  \begin{center}
    \begin{tabular}{lll}
      \toprule
      Symbol & Meaning & Unit \\
      \midrule
      $\x$        & observation point, the head          & m \\
      $\uloc$     & direction a ray {\itshape leaves} $\x$   & -- \\
      $\uext$     & direction a ray {\itshape escapes} the scene & -- \\
      $w$         & power carried by one ray, $1$ at launch & -- \\
      $Q(\uext)$  & share of incident power per steradian of sky & sr$^{-1}$ \\
      $\Phi(\uloc)$ & power arriving per steradian, in units of $S_0$ & sr$^{-1}$ \\
      $\alpha$    & elevation above the horizon          & rad \\
      $\rho$      & slant range from $\x$ to a site      & m \\
      $\nu$       & sites per unit volume                & m$^{-3}$ \\
      \bottomrule
    \end{tabular}
  \end{center}
  \vspace{0.3em}
  Two directions do all the work. $\uloc$ is where a ray starts, $\uext$ is
  where it ends up. They are equal only when nothing was hit.
\end{frame}

\begin{frame}{The move: trace outward, not inward}
  Computing $\chi$ directly means putting a transmitter at every position a base
  station could occupy. Instead everything is computed at the one point where the
  person stands.

  \vspace{0.6em}
  Launch rays from $\x$, follow them until they leave the square, weight each by
  how likely a base station is to sit where it went. Reciprocity licenses this:
  \begin{center}
    a ray leaving along $\uloc$ and escaping along $\uext$ with attenuation $w$
    \\[0.2em]
    $\Longleftrightarrow$
    \\[0.2em]
    a wave from a source at $\uext$ reaching $\x$ from $\uloc$, attenuated by the same $w$.
  \end{center}

  \vspace{0.6em}
  \textbf{One outward trace answers every source position at once.}

  \vspace{0.4em}
  It also buys the rest of the method. Transmitter and receiver sit at the same
  point, so a street level photograph taken from that point sees the surfaces the
  early bounces hit. Facade material is measured, not assumed.
\end{frame}

\begin{frame}{Where the sources are: $Q(\uext)$}
  Base station positions are unknown, so they are a density rather than a list:
  $\nu$ sites per unit volume in a slab, heights $h_{\min}$ to $h_{\max}$ above
  the head, horizontal ranges $r_{\min}$ to $r_{\max}$.

  Count the sites in a narrow cone at elevation $\alpha$:
  \begin{equation*}
    \mathrm{d}N = \nu \dOmega \int_{\rho_-(\alpha)}^{\rho_+(\alpha)} \rho^2 \, \mathrm{d}\rho
    = \frac{\nu}{3}\left[\rho_+^3(\alpha) - \rho_-^3(\alpha)\right] \dOmega ,
  \end{equation*}
  \begin{equation*}
    \rho_-(\alpha) = \max\!\left(\frac{h_{\min}}{\sin\alpha}, \frac{r_{\min}}{\cos\alpha}\right),
    \qquad
    \rho_+(\alpha) = \min\!\left(\frac{h_{\max}}{\sin\alpha}, \frac{r_{\max}}{\cos\alpha}\right).
  \end{equation*}
  Every site is assumed to deliver the same power at the head, because each one
  sets its power to cover its own cell. So the power share follows the count:
  \begin{equation*}
    \boxed{\;Q(\uext) \;\propto\; \rho_+^3(\alpha) - \rho_-^3(\alpha)\;}
    \qquad \text{and} \qquad \oint Q(\uext) \dOmega = 1 .
  \end{equation*}
  Only elevation appears. A uniform density on level ground has nothing to
  distinguish one compass bearing from another.
\end{frame}

\begin{frame}{Two consequences of $Q$}
  \textbf{The support is not a free parameter.} A direction can hold a site only
  where $\rho_- \le \rho_+$, so the height band and the range band already fix
  the elevation limits:
  \begin{equation*}
    \arctan\frac{h_{\min}}{r_{\max}} \;\le\; \alpha \;\le\; \arctan\frac{h_{\max}}{r_{\min}} .
  \end{equation*}

  \vspace{0.5em}
  \textbf{The near horizon carries most of the measure}, and the near horizon is
  where the buildings act. A distant site sits low in the sky, so its path grazes
  along the street and a single facade can intercept it. A site overhead has no
  facade in the way.

  \vspace{0.5em}
  That is why the shape of the square changes the answer at all.

  \vspace{0.5em}
  Normalising $Q$ to unit integral is what makes $\chi$ dimensionless and equal
  to $1$ in free space. The free space denominator $S_0$ is therefore never
  computed.
\end{frame}

\begin{frame}{The estimate}
  A ray launched into $\uloc$ and escaping into $\uext$ collects from the sources
  lying in $\uext$, which carry a share $Q(\uext)$, and delivers a fraction $w$
  of it. Average over the rays launched into $\uloc$, then integrate:
  \begin{equation*}
    \Phi(\uloc) = \big\langle\, w \, Q(\uext) \,\big\rangle_{\uloc} ,
    \qquad
    \boxed{\;\chi = \oint \Phi(\uloc) \dOmega\;}
  \end{equation*}
  In practice the sphere of launch directions is split into $M$ cells of near
  equal solid angle, $N$ rays are launched uniformly at random, $n_c$ of them
  land in cell $c$, and each escaping ray deposits into the cell it was
  \emph{launched} into:
  \begin{equation*}
    \hat{\chi} = \frac{4\pi}{M} \sum_{c=1}^{M} \frac{1}{n_c}
    \sum_{j\,:\,c(j)=c} w_j \, Q(\uext_j) .
  \end{equation*}
  Stratifying by launch cell matters because $Q$ varies by orders of magnitude
  across the sphere. In free space every ray escapes with $w = 1$ and $\uext_j =
  \uloc_j$, so this reduces to a quadrature of $Q$ and returns $1$, as it must.
\end{frame}

\begin{frame}{The split that makes it readable}
  Every ray either escapes untouched or reflects at least once, so the sum
  partitions with nothing left over:
  \begin{equation*}
    \boxed{\;\chi = \chi_{\mathrm{dir}} + \chi_{\mathrm{ref}}\;}
  \end{equation*}
  \begin{itemize}
    \item $\chi_{\mathrm{dir}}$ collects the untouched rays. Under the isotropic
      model it equals the \emph{sky fraction}, the share of directions from the
      head reaching open sky. This is checked as an identity, and it holds to
      $0.0004$\,dB.
    \item $\chi_{\mathrm{ref}}$ is everything the surroundings return, and it is
      strictly positive.
  \end{itemize}
  \vspace{0.5em}
  This separates the only two things a square can do: hide sky, and reflect.
  Both are reported, because a $\chi$ near $1$ can mean an open square or a
  closed one that reflects back what it blocked.
\end{frame}

\begin{frame}{Inside the bounce loop}
  Roughness splits the reflection rather than blurring it. At each interaction
  the fraction of reflected power surviving as a mirror reflection is the
  Rayleigh factor
  \begin{equation*}
    \kappa = \exp\!\left[-\left(\frac{4\pi \sigma_h \cos\theta_i}{\lambda}\right)^{\!2}\right] .
  \end{equation*}
  The outgoing direction is drawn as a mirror reflection with probability
  $\kappa$, and from a cosine weighted lobe about the normal otherwise. Choosing
  one lobe with a probability equal to its share of the energy needs no
  compensating weight, because the two cancel.

  \vspace{0.5em}
  Note the $\cos\theta_i$: a surface that scatters at normal incidence behaves as
  a mirror at grazing incidence. This is the only place that angle dependence
  enters.

  \vspace{0.5em}
  Rays still travelling after $L$ interactions are dropped, and the power they
  carried is reported as a share of the total. The truncation is measured rather
  than assumed.
\end{frame}

\begin{frame}{What the method assumes}
  Three that bound everything else.

  \vspace{0.6em}
  \textbf{Far field.} Every source illuminates the square with a plane wave. The
  closest position any model admits is 10\,m from the head, and $2D^2/\lambda =
  9$\,m at 15\,GHz for the 0.3\,m apertures involved.

  \vspace{0.5em}
  \textbf{Powers add, phases do not.} At 15\,GHz and 100\,m the phase of a path
  turns by about $3.1\times10^{4}$\,rad per radian of exit angle. Resolving it
  would need on the order of $10^{10}$ angular cells against the 512 used here.
  Any phase this method could carry would be noise.

  \vspace{0.5em}
  \textbf{Sites are a density, not a list.} Positions are unknown and are modelled
  as uniform over a slab, bounded in height and range.

  \vspace{0.6em}
  Diffraction is absent by construction. That is the largest known hole and it is
  stated rather than hidden.
\end{frame}

\begin{frame}{The method in one line}
  \begin{center}
    \Large
    Launch rays from the head, weight each by $Q$ where it escapes,
    \\[0.3em]
    and the ratio to free space falls out of one trace.
  \end{center}
  \vspace{1.2em}
  \begin{center}
    \begin{tabular}{ll}
      \toprule
      $\chi = S/S_0$ & the quantity, $1$ in free space \\
      $Q(\uext) \propto \rho_+^3 - \rho_-^3$ & where the sources are \\
      $\Phi(\uloc) = \langle w\,Q(\uext)\rangle_{\uloc}$ & what arrives per steradian \\
      $\chi = \oint \Phi \dOmega$ & the estimate \\
      $\chi = \chi_{\mathrm{dir}} + \chi_{\mathrm{ref}}$ & sky hidden, power returned \\
      \bottomrule
    \end{tabular}
  \end{center}
  \vspace{1.0em}
  \begin{center}
    Co-located transmitter and receiver is what lets a photograph
    \\
    from the standpoint measure the material the first bounce hits.
  \end{center}
\end{frame}

\end{document}
